\setcounter{section}{10}
\section{TỈ SỐ LƯỢNG GIÁC CỦA GÓC NHỌN}
\subsection{Kiến thức trọng tâm}
	\subsubsection{Khái niệm tỉ số lượng giác của góc nhọn}
	\begin{boxdn}
		\immini{
			Cho góc nhọn $\alpha$. Xét tam giác $ABC$ vuông tại $A$ có góc nhọn $B$ bằng $\alpha$. Ta có:
			\begin{itemize}
				\item Tỉ số giữa cạnh đối và cạnh huyền gọi là sin của $\alpha$, kí hiệu $\sin \alpha$.
				\item Tỉ số giữa cạnh kề và cạnh huyền gọi là côsin của $\alpha$, kí hiệu $\cos \alpha$.
				\item Tỉ số giữa cạnh đối và cạnh kề của góc $\alpha$ gọi là tang của $\alpha$, kí hiệu $\tan \alpha$.
				\item Tỉ số giữa cạnh kề và cạnh đối của góc $\alpha$ gọi là côtang của $\alpha$, kí hiệu $\cot \alpha$.
		\end{itemize}}{
			\begin{tikzpicture}[scale=0.65]
				\path 
				(0,0) coordinate (B)
				(3,0) coordinate (A)
				(3,4) coordinate (C);
				\draw (0,0) node[left]{$B$}--(1.5,0) node[below]{cạnh kề}--(3,0) node[right]{$A$}--(3,4) node[above]{$C$}--(0,0);
				\draw (C)--(B)node[midway,sloped,above]{cạnh huyền};
				\draw (A)--(C)node[midway,sloped,below]{cạnh đối};
				\pic[draw, angle radius=3mm]{right angle=B--A--C};
				\pic[draw, angle radius=3mm]{angle=A--B--C};
				\draw (0.3,0.3) node[right]{$\alpha$};
		\end{tikzpicture}}
	\end{boxdn}
	\begin{note}
		\begin{enumEX}[\itemCI]{2}
			\item $\sin \alpha = \dfrac{\text{cạnh đối}}{\text{cạnh huyền}}$;
			\item $\cos \alpha = \dfrac{\text{cạnh kề}}{\text{cạnh huyền}}$;
			\item $\tan \alpha = \dfrac{\text{cạnh đối}}{\text{cạnh kề}}$; 
			\item $\cot \alpha = \dfrac{\text{cạnh kề}}{\text{cạnh đối}}$;
			\item $\cot \alpha = \dfrac{1}{\tan \alpha}$.
			\item! $\sin \alpha$, $\cos \alpha$, $\tan \alpha$, $\cot \alpha$ gọi là các tỉ số lượng giác của góc nhọn $\alpha$.
		\end{enumEX}
	\end{note}
	\begin{note}
		sin, côsin của góc nhọn luôn dương và bé hơn $1$ vì trong tam giác vuông, cạnh huyền dài nhất.
	\end{note}
\begin{vidu}
		Tính các tỉ số lượng giác của góc nhọn $A$ trong mỗi tam giác vuông $ABC$ có $\widehat{B}=90^\circ$ ở hình sau.
	\begin{center}
		\begin{tikzpicture}[declare function={a=2.5;},line join=round,line cap=round,font=\footnotesize,scale=1]
			\path (0:0) coordinate (B)
			(0:a) coordinate (C)
			(B)+(90:0.7*a) coordinate (A)
			($(B)!0.5!(C)$) coordinate (D) node[below,shift={(-90:0.6cm)}] {a)}
			;
			\path pic[draw,angle radius=5pt]{right angle= A--B--C};
			\path pic[draw,angle radius=7pt]{angle= B--A--C};
			\draw (A)--(B) node[midway,left] {$3$}
			--(C) node[midway,below] {$4$}
			--cycle node[midway,right,yshift=3pt] {$5$}
			;
			\foreach \t/\g in {A/90,B/-90,C/-90}{
				\path (\t) node[shift={(\g:7pt)},font=\scriptsize]{$ \t $};
			}
		\end{tikzpicture} \hspace*{3cm}
		\begin{tikzpicture}[declare function={a=2;},line join=round,line cap=round,font=\footnotesize,scale=1]
			\path 
			(0:a) coordinate (C)
			(180:a) coordinate (B)
			(150:a) coordinate (A)
			($(B)!0.5!(C)$) coordinate (D) node[shift={(-90:1cm)}] {b)}
			;
			\path pic[angle eccentricity=2,draw,angle radius=12pt]{angle= A--C--B};
			\path pic[draw,angle radius=5pt]{right angle= B--A--C};
			\draw 
			(A)--(B) node[midway,left] {$1$}
			--(C) node[midway, below] {$\sqrt{17}$}--cycle node[midway, above] {$4$}
			;
			\foreach \t/\g/\i in {A/90/B,B/-90/C,C/-90/A}{
				\path (\t) node[shift={(\g:7pt)}]{$ \i $};
			}
		\end{tikzpicture}
	\end{center}
	\loigiai{
		Xét $\triangle ABC$ vuông tại $B$ (Hình a), ta có
		\begin{enumEX}[\itemCI]{2}
			\item $\sin A=\dfrac{BC}{AC}=\dfrac{4}{5}$;
			\item $\cos A=\dfrac{AB}{AC}=\dfrac{3}{5}$;
			\item $\tan A=\dfrac{BC}{AB}=\dfrac{4}{3}$;
			\item $\cot A=\dfrac{AB}{BC}=\dfrac{3}{4}$.
		\end{enumEX}
		Xét $\triangle ABC$ vuông tại $B$ (Hình b), ta có
		\begin{enumEX}[\itemCI]{2}
			\item $\sin A=\dfrac{BC}{AC}=\dfrac{1}{\sqrt{17}}$;
			\item $\cos A=\dfrac{AB}{AC}=\dfrac{4}{\sqrt{17}}$;
			\item $\tan A=\dfrac{BC}{AB}=\dfrac{1}{4}$;
			\item $\cot A=\dfrac{AB}{BC}=\dfrac{4}{1}$.
		\end{enumEX}
	}
\end{vidu}
	\subsubsection{Tỉ số lượng giác của hai góc phụ nhau}
	\begin{boxdl}
		\immini{
			Nếu hai gọc phụ nhau thì sin góc này bằng côsin góc kia, tang góc này bằng côtang góc kia.
			\begin{note}
				Cho $\alpha$ và $\beta$ là hai góc phụ nhau ($\alpha+\beta=90^\circ$), khi đó
				\begin{itemize}
					\begin{multicols}{2}
						\item $\sin \alpha = \cos \beta$
						\item $\cos \alpha = \sin \beta$
						\item $\tan \alpha = \cot \beta$
						\item $\cot \alpha = \tan \beta$.
					\end{multicols}
				\end{itemize}
		\end{note}}{
			\begin{tikzpicture}[>=stealth,line join=round,line cap=round,font=\footnotesize,scale=0.7]
				\path 
				(0,0) coordinate[label=left:$A$] (A)
				(5,0) coordinate[label=right:$B$] (B);
				\coordinate (D) at ($(B)!1!-30: (A)$) ;
				\coordinate (E) at ($(A)!1!60: (B)$) ;
				\coordinate[label=above:$C$] (C) at (intersection of B--D and A--E);
				\draw (A)--(B)--(C)--cycle;
				\pic[draw, angle radius=2mm]{right angle=B--C--A};
				\pic[draw,double, angle radius=5mm]{angle=C--B--A} node[below] at ($(B) + (145:1.4)$){$\beta$};
				\pic[draw, angle radius=4mm]{angle=B--A--C} node[below] at ($(A) + (35:1.2)$){$\alpha$};
			\end{tikzpicture}
		}
	\end{boxdl}
%%%%%%%%%%%%%%%%%%%
\begin{vidu}
		So sánh
	\begin{enumEX}{2}
		\item $\sin 25^\circ$ và $\cos 65^\circ$;
		\item $ \cos 25^\circ$ và $\sin 65^\circ$;
		\item $\tan 25^\circ$ và $\cot 65^\circ$;
		\item $\cot 25^\circ$ và $\tan 65^\circ$.
	\end{enumEX}
	\loigiai{
		Ta có
		\begin{enumEX}{2}
			\item $\sin 25^\circ=\cos (90^\circ-25^\circ)=\cos 65^\circ$;
			\item $\cos 25^\circ=\sin (90^\circ-25^\circ)=\sin 65^\circ$;
			\item $\tan 25^\circ=\cot (90^\circ-25^\circ)=\cot 65^\circ$;
			\item $\cot 25^\circ=\tan (90^\circ-25^\circ)=\tan 65^\circ$.
		\end{enumEX}
	}
\end{vidu}